%% notes-tex/008-xue-ffa-epistasis/sections/supplementary.tex
%% [[experiments.008-xue-ffa.perspective-epistasis-in-metabolic-engineering]]
%% https://github.com/Mjvolk3/torchcell/tree/main/notes-tex/008-xue-ffa-epistasis/sections/supplementary.tex
%%
%% SOURCE: every number and every panel below is produced by a committed script.
%%   experiments/008-xue-ffa/scripts/supplementary_figure_panels.py
%%       -> notes/assets/images/008-xue-ffa/si_panel_*.svg  (all six SI panels)
%%       -> experiments/008-xue-ffa/results/supplementary_readout_counts.csv
%%          (the per-readout call counts and medians quoted in Notes 2 and 3)
%% The four model equations moved to the Methods section (review 2026.09.16): they define
%% what every figure measures, so they belong with the measurement description rather than
%% behind the main text.

\clearpage
\section*{Supplementary Information\secstatus{todo}}

\noindent\textbf{Contents}\par\smallskip
%% The second column is a wrapping p{} rather than an `l' column, which cannot break, and
%% the 13em reserves what "Supplementary Note~N" plus the column gap actually occupies.
\begin{tabular}{@{}lp{\dimexpr\linewidth-13em\relax}@{}}
\suppnoteref{note:digenic}  & The digenic layer the trigenic score is measured against (with \suppfig{fig:si-digenic})\\
\suppnoteref{note:species}  & The five species under the summed total (with \suppfig{fig:si-species})\\[4pt]
Supplementary Figs.\ & A figure supporting a Note sits with that Note\\
\end{tabular}

%% SI numbering, set ONCE: Supplementary Figures restart at 1 and carry the word
%% "Supplementary" rather than an S prefix, which is the Nature Portfolio form
%% ([[paper.proof-writing-standard]]). There is no equation counter to reset: every
%% equation in this document is now in the main text or in Methods.
\setcounter{figure}{0}
\renewcommand{\figurename}{Supplementary Figure}

\suppnote{note:digenic}{The digenic layer the trigenic score is measured against\secstatus{todo}}

The main text is a claim about the third deletion. That claim is made against what the
first two do, so what the pairs do is worth stating on its own, and it is the opposite of
the trigenic result.

\notesec{The pairs are positive} On total titer, 43 of the 45 double deletions have a
positive multiplicative digenic score and the median is $+0.40$. Twenty-five of the 45
reach a 5\% false discovery rate within the readout, and every one of the 25 is positive
(\suppfig{fig:si-digenic}a). A pair of deletions therefore tends to produce more titer than
the product of its singles, which is the behavior a campaign would want and would find if
it stopped at two.

\notesec{No factor carries it} Drawing all 45 scores as a matrix over the ten factors
(\suppfig{fig:si-digenic}b) shows the positive scores spread across the panel rather than
concentrated in the rows of one or two factors. The largest score in absolute value is
0.98. Every one of the ten factors takes part in at least one called pair, from GCN5 in 1
of its 9 to TFC7 in 8 of its 9, so no single factor carries the positive digenic signal
and none is exempt from it.

\notesec{The third deletion undoes it} Writing the multiplicative expectation for a triple
with the pairs included, as in Eq.~\ref{eq:tau-eps}, gives a prediction of the triple
from its singles and its pairs alone. Ninety-three of the 120 triples fall below that
prediction (\suppfig{fig:si-digenic}c), and the vertical distance to the diagonal in that
panel is $\tau$. The negative trigenic median reported in the main text is therefore not a
continuation of a downward trend the pairs already set. The pairs trend upward, and the
third deletion reverses the direction.

\begin{figure*}[t]
\panelrow{%
  \panel{a}{figures/si_panel_digenic_volcano.pdf}\hfill
  \panel{b}{figures/si_panel_digenic_matrix.pdf}\hfill
  \panel{c}{figures/si_panel_digenic_contribution.pdf}%
}
\caption{\textbf{The digenic layer on total fatty acid titer.} \textbf{a}, Digenic score
$\varepsilon$ against evidence for the 45 double deletions under the multiplicative model.
The dashed line marks the largest $P$ that Benjamini-Hochberg rejects at a 5\% false
discovery rate over the 45 tests of this readout; 25 pairs clear it and all 25 are
positive. \textbf{b}, Every pair's $\varepsilon$ as a matrix over the ten factors, blue
positive and brick negative on a scale centered at zero. The matrix is symmetric and the
dots marking a false-discovery-rate call are drawn in the upper triangle only.
\textbf{c}, Each triple's measured titer against the titer predicted from its singles and
its pairs, $f_i f_j f_k + \varepsilon_{ij} f_k + \varepsilon_{ik} f_j + \varepsilon_{jk}
f_i$. The vertical distance to the diagonal is the trigenic score $\tau$, so a point below
the line is a triple that produces less than its own pairs predict; 93 of 120 do. Filled
points are triples called at a 5\% false discovery rate within total titer.}
\label{fig:si-digenic}
\end{figure*}

\suppnote{note:species}{The five species under the summed total\secstatus{todo}}

Total titer is the sum of five measured fatty acids, and summing is a modeling choice in
the same sense that the null is. The five species were measured separately, so the choice
can be audited rather than argued about.

\notesec{Every species interacts, and they do not agree on the sign} Recomputing the
false-discovery-rate correction within each readout gives 3 to 80 called trigenic
interactions per species against 86 on the summed total
(\suppfig{fig:si-species}a). C16:0, C18:0 and C18:1 behave like the total, with 60 to 67 of
their calls negative. C14:0 does not: all 26 of its calls are positive and its median
trigenic score is $+3.90$, the largest of any readout. C16:1 is nearly silent, with 2
digenic and 3 trigenic calls, so a claim about it would be a claim about a readout this
design cannot resolve.

\notesec{C14:0 reverses between orders} On C14:0 all 15 digenic calls are negative with a
median of $-1.40$, and all 26 trigenic calls are positive with a median of $+3.90$
(\suppfig{fig:si-species}c). The direction of the interaction on that species depends on
how many deletions are counted. Three of the five species move down from the digenic to
the trigenic median and two move up, so the reversal the main text reports on the summed
total is a majority behavior rather than a universal one.

\notesec{The summed total is not a faithful summary} Correlating $\tau$ across readouts
over the same 119 triples (\suppfig{fig:si-species}b) puts the total close to C16:0 and
C18:1 at $r = 0.91$, further from C16:1 and C18:0 at $r = 0.54$ and $0.40$, and on the
other side of zero from C14:0 at $r = -0.28$. The two most negatively related species are
C14:0 and C18:1 at $r = -0.41$. Summing therefore reports the species that carry the most
mass and cancels a species whose interactions run the other way, which is the mechanism
behind the main text's observation that 31 positive consensus interactions exist across the
five species and none survives on their sum.

\notesec{What follows for a campaign} A campaign that optimizes one species is not
optimizing a scaled version of the total, and the interaction structure it would have to
navigate is not the one reported here. The reachability result of the main text is stated
on total titer, and repeating it per species is a separate analysis this document does not
carry out.

\begin{figure*}[t]
\panelrow{%
  \panel{a}{figures/si_panel_species_counts.pdf}\hfill
  \panel{b}{figures/si_panel_species_agreement.pdf}\hfill
  \panel{c}{figures/si_panel_species_medians.pdf}%
}
\caption{\textbf{The five measured species behind the summed total.} \textbf{a},
Interactions called at a 5\% false discovery rate within each readout, digenic (solid) and
trigenic (hatched), stacked by sign, blue positive and brick negative. Counts above each
bar are the totals. \textbf{b}, Trigenic $\tau$ compared across readouts over the same 119
triples: Pearson $r$ below the diagonal, and above it the number of triples called at a 5\%
false discovery rate in both readouts. \textbf{c}, Median interaction score per readout for
both orders. The gap between the two lines is the quantity the main text reports on total
titer, and its sign is not the same on every species.}
\label{fig:si-species}
\end{figure*}
